\chapter{Background and Related Work}

\section{Visual question answering and GQA}
\label{sec:ch2-vqa}

Visual question answering (VQA) maps an image and a natural-language question to an answer. The output may be generated as text or selected from a fixed answer vocabulary; the latter turns VQA into supervised classification and makes the answer space part of the task definition \citep{antol2015vqa}. A closed vocabulary offers a stable comparison target, but it also limits coverage and makes results conditional on which answers are included. Early work showed that comparatively direct combinations of image and question features can form competitive VQA baselines, including concatenation followed by a classifier \citep{zhou2015ibowimg,jabri2016revisiting}. These systems establish an important reference point: sophisticated multimodal processing should be compared with well-defined simple alternatives rather than treated as necessary by default.

GQA was designed around questions derived from scene graphs and provides functional programs associated with those questions \citep{hudson2019gqa}. The program annotations make it possible to group examples by the length of the recorded operation sequence. Program length is a proxy for reasoning demand, not a ground-truth measure of internal reasoning depth: different programs can vary in linguistic form, visual content and answer distribution even when their lengths match. It is therefore useful for stratified evaluation, but it does not isolate reasoning from every other source of difficulty.

Answer frequencies and language regularities can also influence behaviour on GQA. Work on out-of-distribution and reasoning-oriented evaluations shows that changes in answer priors and statistical cues affect what a model can predict without establishing the intended visual relation \citep{kervadec2021gqaood,kervadec2021reasoning}. Closed-answer evaluation must consequently distinguish evidence of task performance from evidence about the process that produced it. This boundary motivates the dissertation's separate treatment of unimodal controls, visual interventions and compositional profiles while leaving the formal research questions in Section~\ref{sec:ch1-rqs}.

\section{Frozen vision--language representations}
\label{sec:ch2-representations}

Contrastive vision--language pretraining can place images and text in a shared representation space, as established by the foundational CLIP approach. % CITE-GAP clip_foundation
SigLIP is a related representation-learning approach built around a sigmoid loss rather than the batch-normalised contrastive objective. % CITE-GAP siglip
These foundational accounts describe representation-learning concepts, not the exact weights used in this dissertation. The implementation checkpoints and their provenance are specified separately in Chapter~3. In the headline systems studied here, the visual and textual encoders remain frozen: training changes only the downstream interface or head. This separation allows downstream architectures to be compared without simultaneously adapting the representation source.

Freezing an encoder turns its representations into a fixed transfer substrate. This reduces the trainable scope, but it also means that a downstream model can use only information exposed by the retained features and accessible to its architecture. Linear or small nonlinear classifiers over frozen CLIP embeddings provide precedents in closed-answer visual tasks, while larger fusion networks show that substantial downstream processing can still be placed above the same broad feature type \citep{deuser2022less,bourigault2025frevl}. Their tasks, head capacities and protocols differ, so this relationship is architectural rather than a numerical comparison.

The choice of representation granularity is separate from the choice of downstream head. A global vector offers a compact summary, whereas token-level features retain spatially distributed inputs that a later module may attend over. Classification studies have found closely related outcomes for several pooling choices in some vision-transformer regimes, while attentive probes can expose information not available to a linear readout, with the benefit depending on the backbone \citep{zhai2022scaling,psomas2026ep}. These observations motivate treating representation quality, feature granularity and downstream capacity as distinct design axes rather than assuming that more tokens or a larger head must improve a task.

\section{Multimodal fusion for lightweight heads}
\label{sec:ch2-fusion}

The simplest multimodal combination concatenates image and question vectors and lets a classifier learn any useful dependence. Explicit interaction templates instead expose relations before the learned head. The template $[i;q;\lvert i-q\rvert;i\odot q]$, originating in sentence-pair modelling, combines the two inputs with an absolute difference and elementwise product \citep{conneau2017infersent}. Applied to aligned vision--language embeddings, these terms offer fixed comparison and agreement features without requiring a learned attention module.

Bilinear and factorised fusion methods provide a more general family of multiplicative interactions. Their designs trade representational flexibility against the number and organisation of trainable parameters; MUTAN, for example, frames fusion architecture and parameter budget together rather than treating them as independent \citep{benyounes2017mutan}. Attention-based fusion introduces another form of interaction by allowing one modality to weight elements of the other. The structure chosen for fusion determines which cross-modal relations are explicit and which must be learned indirectly by later layers. None of these constructions is universally preferable: the useful inductive bias depends on the available representation, data regime, output task and capacity assigned to the comparison.

This dependence creates a control problem. If an interaction model has a wider input or larger hidden layer than concatenation, a difference between them may reflect capacity as well as feature structure. Related work has compared concatenation and cross-attention across data scales on image--text matching, while EMAP tests a trained multimodal model against an additive approximation to diagnose whether cross-modal interactions materially affect its predictions \citep{zhou2026alignment,hessel2020emap}. These are different tasks and diagnostics, but they motivate the capacity controls and explicit interaction ablations described in Chapter~3 without predetermining their VQA outcome.

\section{Latent-query architectures and connectors}
\label{sec:ch2-latent}

Latent-query architectures replace direct processing of every input pair with a compact learned bottleneck. The Perceiver uses a fixed set of latent vectors that cross-attend to a larger input and then interact internally, separating input size from the main processing depth \citep{jaegle2021perceiver}. Flamingo's Perceiver Resampler applies this principle as a connector that converts variable visual features into a fixed set of tokens for a language model \citep{alayrac2022flamingo}. In both cases, learned queries organise access to a larger frozen or pretrained representation rather than serving as answer classes themselves.

BLIP-2 introduces a Q-Former between a frozen image encoder and a language model. Its learned query embeddings extract visual information for the language-model interface, and the Q-Former is pretrained as a connector before downstream use \citep{li2023blip2}. It uses 32 learned queries. MAPL likewise uses learnable query tokens in a comparatively small mapping network between frozen components \citep{manas2023mapl}. The common architectural idea is selective access through learned queries; the surrounding training objective and destination of those queries remain central to what the module represents.

Connector comparisons caution against reading architectural sophistication as an outcome by itself. In its frozen-encoder and frozen-language-model regime, DePALM reports that a simpler global query-pooling mapper can compare favourably with local resampler variants, particularly under limited training data. MM1 places connector choice below image resolution and visual-token count in comparative importance, whereas IDEFICS2 provides a counterpoint in which learned latent pooling improves its own configuration \citep{vallaeys2024depalm,mckinzie2024mm1,laurencon2024idefics2}. These findings come from generative systems with different scales and training histories, so they support regime-sensitive comparison rather than a general ranking of connectors.

The dissertation's reasoner follows the broader learned-query interaction mechanism and also uses 32 latents, but it is not a Q-Former. It is a discriminative reasoning head over frozen question and visual tokens for closed-set VQA, not a connector that prepares visual tokens for a generative language model. It is trained on the task data rather than through large-scale connector pretraining, and it differs in role, scale and output objective. The shared query count therefore does not establish functional equivalence, and no performance comparison with Q-Former systems follows from the architectural lineage.

\section{Frozen language models as multimodal interfaces}
\label{sec:ch2-lm-interfaces}

Frozen language models can receive non-text information through a learned interface. Frozen trained a visual prefix for a fixed language model, LiMBeR used a single linear map from frozen image features into a frozen language model, and eP-ALM used a restricted projection and soft-token interface between frozen components \citep{tsimpoukelli2021frozen,merullo2023limber,shukor2023epalm}. These approaches differ in how the visual component is obtained and how the output is evaluated, but they show that multimodal adaptation need not require updating the language model itself.

Placement changes the function assigned to a language model. It can encode a question before multimodal interaction, supplying contextual token features to a separate reasoner, or it can consume a multimodal state after interaction and participate in the answer readout. The dissertation uses frozen SmolLM2 models to study these two interfaces. % CITE-GAP smollm2
For each placement, an architecture-matched random-initialised frozen model provides a control for pretraining while holding the model configuration fixed. This is a controlled design choice of the dissertation, not a claim that earlier interfaces are inadequate or absent.

The interface distinction also prevents model-family size from standing in for a causal account of pretraining. A pretrained model and a random control must be compared at the same architecture and placement, while comparisons across placements remain system-level because the information supplied to and requested from the language model differs. Such a matched control addresses the contribution of pretraining within an interface; it does not by itself identify an optimal placement. Chapter~3 defines those interfaces and controls precisely. Chapter 2 uses the prior connector literature only as architectural context and makes no claim about the dissertation's measured outcomes.

\section{Compositional reasoning and representation limits}
\label{sec:ch2-composition}

GQA's functional programs make multi-step grouping possible, but difficulty across program lengths is affected by more than a single reasoning operation. Hop-stratified evaluation of multimodal language models reports declining performance as question chains become longer, while GQA studies show that answer distributions and language priors can influence predictions \citep{kil2024iimmr,kervadec2021gqaood,kervadec2021reasoning}. A program-length profile is therefore best read as a prior-sensitive task diagnostic. It can reveal that performance changes across question structures, but cannot by itself identify an internal reasoning process or localise the cause of a deficit.

Compositional benchmarks and probes examine whether frozen representations preserve relations, order and binding. ARO and Winoground construct controlled text--image contrasts that expose sensitivity to relational structure, while probe work on frozen CLIP embeddings reports difficulty with binding relations \citep{yuksekgonul2023aro,thrush2022winoground,lewis2024bind}. A counterpoint is that linear probes can recover attribute--object binding information inside unimodal embeddings even when the cross-modal alignment makes that information harder to use \citep{koishigarina2026bow}. Together, these results distinguish information being absent from information being present but inaccessible to a particular comparison or interface.

The literature also offers competing accounts of where multimodal limitations arise. MMVP traces some failures of multimodal language models to blind spots inherited from the visual encoder, whereas Fu locates failures in how a downstream language model uses visual features and Takishita shows that large decoders can compensate for degraded visual contextualisation in some regimes \citep{tong2024mmvp,fu2025hidden,takishita2025compensate}. These positions need not be mutually exclusive: the effective limitation can depend on the encoder, downstream capacity, training regime and task. They motivate controlled localisation, but none establishes the mechanism behind this dissertation's observations.

\section{Efficiency in compact multimodal systems}
\label{sec:ch2-efficiency}

Efficiency is not a single interchangeable quantity. Trainable parameter count describes the capacity updated for a task, while total loaded parameters describe a broader resident system; neither is a runtime measurement. Systems with similar parameter counts can execute different operations and move different quantities of intermediate state. Partial-pipeline timing excludes some stages that end-to-end latency includes, and even two end-to-end measurements are comparable only when hardware, precision, batching, preprocessing and timing boundaries align. Fusion and connector studies variously emphasise parameter budgets, token reduction or task performance under compact interfaces rather than one common runtime contract \citep{benyounes2017mutan,vallaeys2024depalm,hu2024mqt}.

Compact instruction-tuned vision--language models provide a useful contextual system class, including SmolVLM. % CITE-GAP smolvlm
They are integrated pretrained models rather than project-developed heads over the dissertation's frozen representations. Their scale labels, trainable parameter counts and reported timings answer different questions unless the measurement boundary is explicitly matched. Chapter~3 therefore defines the project's capacity quantities separately from runtime, and Chapter~7 defines its valid end-to-end warm-serial protocol. Chapter 2 motivates that separation and makes no direct cross-paper latency comparison.

\section{Positioning of this dissertation}
\label{sec:ch2-positioning}

The literature reviewed above addresses different subsets of the relevant design choices. It provides frozen vision--language representations, lightweight classifiers over global features, explicit and learned fusion mechanisms, latent-query connectors, frozen-language-model interfaces, compositional probes and several forms of compactness. It also shows why conclusions about representation, interaction and downstream processing must remain tied to the task and training regime in which they were obtained.

Against this background, the dissertation studies representation quality, explicit multimodal interaction, lightweight latent reasoning and pretrained-versus-random small-language-model interfaces within one controlled closed-set VQA framework. It separates matched capacity from interaction structure, treats question-side and answer-side language-model placement as distinct interfaces, and measures end-to-end runtime rather than inferring it from parameter counts. The chapters that follow specify the governed methods and evidence; this positioning does not claim priority or use the project's results to adjudicate the broader literature.
